\section*{Chapitre \ref{ch:b1:electrophilic-additions} --- Alcènes et alcynes~: les additions électrophiles}

\begin{solution}{exo:b1:electrophilic-additions:1}
2-Chloropropane~; 2-bromo-2-méthylpropane~; 1-bromo-1-méthylcyclohexane~;
2-chlorobut-2-ène.
\end{solution}

\begin{solution}{exo:b1:electrophilic-additions:2}
1,2-Dibromoéthane~; \textit{trans}-1,2-dibromocyclohexane~;
1,2-dibromopropane.
\end{solution}

\begin{solution}{exo:b1:electrophilic-additions:3}
Propan-2-ol~; 2-méthylpropan-2-ol~; 1-méthylcyclohexan-1-ol (Markovnikov
dans chaque cas).
\end{solution}

\begin{solution}{exo:b1:electrophilic-additions:4}
Le propyne, \omterm{def:b1:electrophilic-additions:acetylide}{alcyne terminal}, réagit~: \ce{CH3C#CH + NH2- -> CH3C#C- + NH3}.
L'éthanol~:
\ce{CH3CH2OH + NH2- -> CH3CH2O- + NH3}. Le but-2-yne et le propène n'ont
pas d'hydrogène assez acide.
\end{solution}

\begin{solution}{exo:b1:electrophilic-additions:5}
Le doublet $\pi$ capte \ce{H+} sur le carbone du \ce{CH2}~: cation
tertiaire \ce{(CH3)3C+}~; un \omterm{def:b1:lewis-resonance-vsepr:lewis}{doublet non liant} de l'eau s'y lie~:
\ce{(CH3)3COH2+}~; une molécule d'eau capte un proton~: \ce{(CH3)3COH}.
Dans l'acide concentré chaud, l'eau est rare et l'alcène volatil
s'échappe~: la réaction inverse, la \omterm{def:b1:alcohol-activation:dehydration}{déshydratation}, a lieu.
\end{solution}

\begin{solution}{exo:b1:electrophilic-additions:6}
La protonation du 2-méthylpropène donne un cation tertiaire, celle du
propène un cation secondaire. L'\omterm{def:b1:elementary-steps:energy-profile}{état de transition} de la protonation,
\omterm{def:b1:elementary-steps:rds}{étape cinétiquement déterminante}, ressemble au cation (Hammond)~: le
cation tertiaire, plus stable, est atteint en franchissant une barrière
plus basse.
\end{solution}

\begin{solution}{exo:b1:electrophilic-additions:7}
L'ion bromonium ponte une face du cycle~; l'ion bromure l'ouvre par
l'autre face~: les deux bromes se retrouvent en \textit{trans}. Le
dérivé dibromé \textit{trans} est \omterm{def:b1:stereochemistry-in-depth:stereocentre}{chiral} (pas de plan de symétrie), mais
l'ion bromure attaque également les deux carbones~: les deux
\omterm{def:b1:stereochemistry-in-depth:enantiomers}{énantiomères} se forment en quantités égales, et le produit est racémique,
optiquement inactif.
\end{solution}

\begin{solution}{exo:b1:electrophilic-additions:8}
La protonation de C1 donne le cation secondaire sur C2, voisin de C3,
tertiaire et porteur d'un hydrogène~; cet hydrogène migre avec son
doublet vers C2 (une migration-1,2 d'hydrure), laissant un cation
tertiaire sur C3, capturé par
\ce{Cl-}~: 2-chloro-2-méthylbutane.
\end{solution}

\begin{solution}{exo:b1:electrophilic-additions:9}
Le propyne et l'amidure de sodium donnent l'ion propynure~; avec le
1-bromopropane, par SN2, on obtient
\ce{CH3C#C-CH2CH2CH3}, l'hex-2-yne. Le propynure avec le benzaldéhyde,
puis un acide dilué~:
\ce{C6H5CH(OH)C#CCH3}, le 1-phénylbut-2-yn-1-ol.
\end{solution}

\begin{solution}{exo:b1:electrophilic-additions:10}
\cip{E}-but-2-ène~: ion bromonium sur une face, ion bromure par l'autre
sur C2 ou sur C3~; les deux attaques donnent le même composé, de
descripteurs \cip{2R,3S}~:
le \omterm{def:b1:stereochemistry-in-depth:enantiomers}{composé méso}, un seul stéréoisomère. \cip{Z}-but-2-ène~: l'attaque
sur C2 donne \cip{2S,3S} (ou son image dans un miroir à partir de
l'autre face), l'attaque sur C3
\cip{2R,3R}~; en quantités égales~: deux \omterm{def:b1:stereochemistry-in-depth:isomers}{stéréoisomères}, un \omterm{def:b1:stereochemistry-in-depth:enantiomers}{mélange
racémique}.
\end{solution}

\begin{solution}{exo:b1:electrophilic-additions:11}
Dans l'ion bromonium dissymétrique, le carbone le plus substitué porte
une plus grande part de la charge positive (sa liaison C--Br est plus
longue et plus faible, comme dans un cation presque tertiaire ou
secondaire). L'eau, en large excès, attaque ce carbone par la face
opposée au pont~: 1-bromopropan-2-ol.
\end{solution}

\begin{solution}{exo:b1:electrophilic-additions:12}
$D = 1$~: un alcène. Le 2-méthylbut-1-ène et le 2-méthylbut-2-ène donnent
tous deux le cation tertiaire, d'où le 2-bromo-2-méthylbutane et le
2-méthylbutan-2-ol. Leurs spectres de RMN du proton diffèrent~: deux
protons vinyliques (\ce{=CH2}, singulets) pour le premier, un seul
(\ce{=CH-}, un quadruplet) pour le second.
\end{solution}

\begin{solution}{pb:b1:electrophilic-additions:1}
\textbf{1.} Sur chaque carbone, \ce{CH3} est prioritaire sur H~: le
\cip{Z} a les deux méthyles du même côté, l'\cip{E} de part et d'autre.
\textbf{2.} Des \omterm{def:b1:stereochemistry-in-depth:isomers}{stéréoisomères} qui ne sont pas images l'un de l'autre
dans un miroir~: des \omterm{def:b1:stereochemistry-in-depth:enantiomers}{diastéréoisomères}.
\textbf{3.} La rotation autour de C=C romprait la liaison $\pi$, ce qui
coûte bien plus d'énergie que n'en fournissent les chocs à température
ambiante.
\textbf{4.} L'\cip{E}~: ses groupes méthyle ne se gênent pas.
\textbf{5.} \ce{C4H8 + Br2 -> C4H8Br2}.
\textbf{6.} Deux \omterm{def:b1:stereochemistry-in-depth:stereocentre}{centres stéréogènes} équivalents~: au plus quatre
\omterm{def:b1:stereochemistry-in-depth:isomers}{stéréoisomères}, en réalité trois
(\cip{2R,3R}, \cip{2S,3S} et le méso \cip{2R,3S}).
\textbf{7.} Le doublet $\pi$ attaque un atome Br de \ce{Br2}, le doublet
Br--Br part sous forme de \ce{Br-}, et un \omterm{def:b1:lewis-resonance-vsepr:lewis}{doublet non liant} du brome
attaqué se lie au second carbone~: un ion bromonium à trois atomes.
\textbf{8.} Dans l'ion ponté, chaque atome a un octet~; un \omterm{def:b1:electronic-effects:intermediates}{carbocation}
secondaire ouvert aurait un carbone à six électrons.
\textbf{9.} Un \omterm{def:b1:lewis-resonance-vsepr:lewis}{doublet non liant} de \ce{Br-} se lie à un carbone du
cycle~; la liaison C--Br du pont se rompt, son doublet passant sur le
brome pontant.
\textbf{10.} Le pont occupe une face~; comme dans la SN2, le \omterm{def:b1:electronic-effects:nucleophile}{nucléophile}
entre à l'opposé de la liaison qui se rompt.
\textbf{11.} Une \omterm{def:b1:electrophilic-additions:halonium}{addition anti}.
\textbf{12.} L'eau concurrencerait l'ion bromure auprès de l'ion
bromonium et donnerait une bromhydrine.
\textbf{13.} \cip{2R,3S} (ou \cip{2S,3R}, le même composé).
\textbf{14.} Le même composé.
\textbf{15.} Dans une \omterm{def:b1:stereochemistry-in-depth:isomers}{conformation} où les deux bromes sont éclipsés, un
plan de symétrie passe entre C2 et C3~; les descripteurs sont opposés.
\textbf{16.} \cip{2R,3R} et \cip{2S,3S}.
\textbf{17.} En quantités égales~: un \omterm{def:b1:stereochemistry-in-depth:enantiomers}{mélange racémique}.
\textbf{18.} Oui~: les deux alcènes \omterm{def:b1:stereochemistry-in-depth:enantiomers}{diastéréoisomères} donnent des
produits différents (méso contre racémique).
\textbf{19.} Aucun~: les deux \omterm{def:b1:stereochemistry-in-depth:enantiomers}{énantiomères} se compensent.
\textbf{20.} $5.60/56.0 = \qty{0.100}{mol}$ de chacun.
\textbf{21.} $0.100 \times 0.85 \times 215.8 = \qty{18.3}{g}$ pour chacun.
\textbf{22.} \textbf{$[\alpha] = 0^\circ$ (\omterm{def:b1:stereochemistry-in-depth:enantiomers}{composé méso}), \qty{18.3}{g}}.
\end{solution}
